Prompt templates often differ in small ways that look harmless: a trailing space, a template written in capitals, missing punctuation, or an input wrapped in JSON. We ask whether such meaning-preserving edits change what instruction-tuned language models output, and whether larger models are less affected. We apply single, atomic surface edits to chat-formatted prompts and run ten open-weight models (\qwen 0.5B--14B, \llama 1B/8B, \gemma 9B) with greedy decoding on about 235k generations. We measure classification accuracy, response length, instruction compliance on \ifeval, and extraction F1. A replicate run and placebo edits measure the noise floor of the inference stack. Across 40 model$\times$task cells, accuracy spread over ten variants has a median of 5.8 points (range 1.0--28) and exceeds a permutation null in 27 cells. Sensitivity falls with scale ($-9.5$ points of spread per tenfold increase in parameters, $\rho=-0.87$). The decline is not log-linear, though: it plateaus at about 4 points between 7B and 14B. ALL-CAPS templates are the most harmful edit for small models (up to $-27$ points) and still lower \ifeval compliance by 3.1--4.4 points at 7--14B. Two common beliefs do not hold: trailing whitespace does not lengthen outputs (length ratios 0.94--1.07), and JSON-structured prompts do not give a uniform 10--15 point extraction gain (effects range from $-17$ to $+26$ F1 depending on task and model). Which edits help or hurt varies by model, even between the 16-bit and 4-bit versions of the same weights. We recommend evaluating small models over several formats and avoiding ALL-CAPS templates at any size.
